\chapter{基于类别解耦与动态路由的分级安全知识编辑方法}\label{chap:graded-edit}

本章针对多类别安全编辑中的交叉干扰与过度拒答，研究基于类别解耦与动态路由的分级安全知识编辑方法。独立低秩模块共享模型计算路径。若路由将请求判入错误类别，所调用的模块仍可能改变该请求的响应。为限制此类影响，训练时约束各模块在其他类别和无害输入上的局部输出增量。推理时由预测类别选择模块，由预测风险等级确定调用强度，使同一类别的不同请求可以接受不同程度的干预。

本章分两步检验这些设计。先在固定调用条件下测量模块对非目标类别和无害请求的影响，再在预测路由下比较固定强度与分级强度的实际处置效果。算法实现、风险标注与实验结果均待完成。

\section{相关知识介绍}
\subsection{大模型过度拒答问题}

过度拒答指模型对符合使用策略的正常请求给出不必要的拒绝。XSTest构造了措辞相近的安全和不安全提示，用于检验模型是否因敏感用语而拒绝正常请求\cite{rottger2024xstest}。OR-Bench以表面上可能触发拒绝、实际无害的提示扩大评估范围\cite{cui2025orbench}。因此，本文将相邻无害请求单列，并与一般无害请求、风险请求同时评估。本文沿用第二章的ASR、SRR与效用指标，并定义无害误拒率
\begin{equation}
 \mathrm{ORR}=\frac{1}{|\mathcal T_u|}\sum_{x\in\mathcal T_u}
 \mathbb I[\text{对正常请求}x\text{作出不必要拒绝}].
 \label{eq:ch3-orr}
\end{equation}
是否构成不必要拒绝，应依据统一规则和人工抽样复核判断；不能只检索固定拒答关键词，也不能把一切拒答都作为安全成功。

\subsection{多类别分级安全编辑}

在一般事实编辑中，Gupta等观察到连续编辑会影响旧编辑的保持与模型通用能力\cite{gupta2024forgetting}。另一方面，RippleEdits指出，编辑一个事实后，相关事实的回答也可能需要随之改变\cite{cohen2024ripple}。上述结果说明，编辑后的非目标变化需要结合任务判断。本文分别记录非目标类别的正、负变化，利用逐类干扰矩阵检验不同风险类别的实际影响。

设风险类别集合为$\mathcal C=\{1,\ldots,C\}$，原始模型为$f_{\theta_0}$。本章以同一原始模型为起点，为各类别学习独立参数增量$\Delta W_c$，推理时按输入选择模块。独立存储避免了对同一份基座权重反复覆盖，但各模块仍作用于共享计算路径，因此不自动保证行为互不干扰。

主实验拟先选择SafeEdit中样本充分、语义有区分的四类风险，补充相邻无害样本。类别名称、样本量和选择理由待数据审计后填写；不假定已完成标注。风险等级通过路由强度表达，每类只训练一个低秩模块，不为每个“类别×等级”额外训练一套模型。

\subsection{风险类别与等级划分}

安全基准使用的风险标签在命名和粒度上存在差异，Rastegar等通过层次化分类对齐不同基准的标签\cite{rastegar2026taxonomy}。风险类别还可以与独立的等级维度搭配。例如，CRiskEval按类别和安全等级评估中文场景中的前沿风险倾向\cite{shi2025criskeval}。本文的等级用于标注生成响应所需的处置方式，由当前任务的规则和样本语境确定。

类别表示请求所属的语义风险，等级表示所需的响应方式。能够直接回答的正常请求记为低风险。需要在安全边界内提供说明或替代信息的请求记为中风险。需要拒绝危险部分并给出适当安全建议的请求记为高风险。标注时先判断请求意图、所求输出和使用情境，再确定能提供的信息范围。中、高风险使用同一安全边界。含有风险词汇的正常科普或预防请求仍按其实际意图标注。

训练样本拟记录$(x,c,r,y^{\mathrm{safe}})$。低风险样本的类别字段可留空，类别损失只在具有确定标签的中、高风险样本上计算。对多标签、歧义和类别边界样本另设标记，并在评估中单列。风险等级须按上述规则另行标注，SafeEdit原有标签用于确定候选类别\cite{wang2024safeedit}。由两名标注者独立复核抽样数据，记录分歧及裁定依据。等级标注及一致性统计\textbf{待完成}。

\subsection{低秩编辑与动态路由}

LoRA为以少量参数表达模型更新提供基础，MELO进一步组织可动态激活的低秩模块，ToxEdit探索根据毒性激活调整计算路径\cite{hu2022lora,yu2024melo,lu2025toxedit}。相关方法说明模块化表示与输入自适应干预具有可行的技术基础。本文围绕风险类别与处置等级，进一步研究模块作用范围约束和类别级调度的协同关系。

对于连续知识编辑，KDE保留先前编辑的表征基并对新梯度作正交投影，WISE在原有知识与编辑记忆之间按输入路由，并将不同批次编辑分配到参数子空间\cite{xu2025kde,wang2024wise}。本章以风险类别作为低秩模块的组织单位，通过非目标输入上的功能约束限制类别外影响。在按需干预方面，SafeSwitch利用内部激活判别风险并选择是否启用安全分支\cite{han2025safeswitch}。本文的路由还根据预测类别选择模块，并按处置等级调整调用强度。

本章以类别模块为独立更新单元，以非目标功能约束控制单个模块的作用范围，以风险分级门控调节实际施加的更新幅度。实验设置共享低秩编辑、无解耦约束和固定强度调用三类对照，分别考察模块拆分、作用范围约束及调用强度的影响。

\section{方案框架}
\subsection{多类别风险识别与等级划分}

整体流程分为离线训练和在线推理。离线阶段先根据训练数据上的安全/不安全响应表示差异确定局部编辑位置，该定位环节采用既有安全编辑机制\cite{wang2024safeedit}。主方案汇总各类别训练样本的层差异分数，固定一个共享编辑层及其MLP输出投影，具体聚合规则在实现后记录。随后构造类别模块，以非目标功能约束训练，并在独立验证集校准风险识别和路由阈值。

在线阶段提取冻结原始模型在编辑层之前的提示表示$h(x)$，以轻量分类头输出三级概率$q_r(x)=p(r\mid x)$及条件类别分布$p(c\mid x,\mathrm{risk})$，其中$r\in\{0,1,2\}$。定义归一化处置分数$s(x)=q_1(x)/2+q_2(x)\in[0,1]$。特征只使用提示，不使用待生成回答或测试标签；同一输入只计算一次路由，生成期间固定，避免路由与已编辑激活相互反馈。拟采用线性头或小型MLP，不额外训练大型检测模型；数据划分按提示族去重，并分别记录用途。

\begin{figure}[htbp]
 \centering
 \fbox{\begin{minipage}{0.88\textwidth}\centering
 离线：局部编辑位置确定 $\rightarrow$ 类别低秩模块 $\rightarrow$ 解耦训练\\[3pt]
 提示及标注 $\rightarrow$ 冻结表示分类头 $\rightarrow$ 验证集校准\\[6pt]
 在线：输入 $\rightarrow$ 风险及类别识别 $\rightarrow$ 分级门控\\[3pt]
 $\rightarrow$ 旁路或类别模块 $\rightarrow$ 响应生成 $\rightarrow$ 独立评估
 \end{minipage}}
 \caption{多类别分级解耦与动态路由方案设计示意（待实现）}
 \label{fig:ch3-framework}
\end{figure}

\subsection{基于类别独立低秩模块的安全知识编辑}

对选定权重$W\in\mathbb R^{d_{\mathrm{out}}\times d_{\mathrm{in}}}$，定义
\begin{equation}
 \Delta W_c=\frac{\gamma}{k}B_cA_c,\qquad
 A_c\in\mathbb R^{k\times d_{\mathrm{in}}},\quad
 B_c\in\mathbb R^{d_{\mathrm{out}}\times k},
 \label{eq:ch3-lora}
\end{equation}
其中$k$为低秩维度，$\gamma$为缩放常数。基座冻结，训练第$c$类模块时不更新其他模块，初始化采用一侧随机、一侧为零的LoRA方式。类别编辑目标沿用式\eqref{eq:safety-edit-loss}，在本类别安全响应样本上计算，并保留无害输出KL约束。训练数据和计算预算应与共享LoRA对照保持可比。

所有类别模块使用相同编辑层和秩配置。单层模块参数总量为$Ck(d_{\mathrm{in}}+d_{\mathrm{out}})$，不含分类头；单次Top-1调用只计算一个低秩分支。模块独立存储不等于显存自动减少，具体加载策略及开销需实测。

\subsection{类别间参数干扰分析与解耦约束}

先进行DINM逐类编辑预实验，在验证子集上观察只编辑第$i$类时其他类别及无害输入的响应变化。正式评估在独立测试集上构造干扰矩阵
\begin{equation}
 M_{ij}=\mathrm{SRR}_j(f_i)-\mathrm{SRR}_j(f_{\theta_0}),
 \label{eq:ch3-interference}
\end{equation}
其中$f_i$为从同一基座出发且仅施加第$i$类编辑的模型，$\mathrm{SRR}_j$在第$j$类独立测试集上计算。对角项衡量目标类变化。非对角正值可能是有益迁移，负值表示安全退化。各行另记录无害效用与ORR。诊断类别模块时，对每条测试输入固定以单位强度调用第$i$类模块，关闭预测路由、低风险旁路及安全回退。这样得到的是同一模块作用于目标和非目标输入时的影响。使用预测路由的效果在端到端评估中另行报告。预实验及矩阵数值\textbf{待完成}。

类别参数分离不保证功能分离，也不要求语义重叠的类别严格正交。主方案约束模块在非目标输入上的作用。设$\mathcal D_{-c}$包含其他类别及无害训练样本，$H(x)$为冻结基座在所编辑投影入口处的表示，定义
\begin{equation}
 \mathcal R_{\mathrm{dec}}
 =\frac1C\sum_{c=1}^{C}\frac1{|\mathcal D_{-c}|}
 \sum_{x\in\mathcal D_{-c}}
 \frac{\|\Delta W_cH(x)\|_2^2}{d_{\mathrm{out}}}.
 \label{eq:ch3-decouple}
\end{equation}
多token时对有效提示位置取均值。第$c$类模块被调用时，其局部输出增量由$\Delta W_c H(x)$按调用强度缩放。对非目标输入约束该增量，旨在减小误选模块时可能施加的扰动。该项直接约束所选层的局部输出，最终生成行为仍须实测。由于约束也可能抑制有益迁移，权重由验证集选择，评估同时报告目标类安全、非目标类变化和无害误拒率。

\subsection{基于风险感知的动态路由权重计算}

设$0<\tau_1<\tau_2<1$为两个分级阈值，$0<\beta<1$为预设的中风险模块强度。定义
\begin{equation}
 a(x)=\begin{cases}
 0,&s(x)<\tau_1,\\
 \beta,&\tau_1\le s(x)<\tau_2,\\
 1,&s(x)\ge\tau_2.
 \end{cases}
 \label{eq:ch3-risk-gate}
\end{equation}
主方案采用Top-1，$c^*(x)=\arg\max_c p(c\mid x,\mathrm{risk})$，最终权重为
\begin{equation}
 w_c(x)=a(x)\mathbb I[c=c^*(x)],\qquad
 W(x)=W+\sum_{c=1}^{C}w_c(x)\Delta W_c.
 \label{eq:ch3-route}
\end{equation}
式\eqref{eq:ch3-route}中，$c^*(x)$确定调用哪个类别模块，$a(x)$确定施加该模块增量的比例。预测为低风险时旁路类别模块，预测为中、高风险时分别按$\beta$和$1$调用。真实处置等级由标注给定，用于评估预测错误及其对响应的影响。Top-2加权组合仅作为多标签场景的可选扩展。

分类头使用三级监督交叉熵及仅在风险样本上计算的类别交叉熵。二分类风险概率取$1-q_0(x)$，与处置分数$s(x)$分别记录。分级阈值在等级验证集上校准，优先满足预先设定的高风险召回要求，再降低低风险误触发。高风险门控下若$\max_c p(c\mid x,\mathrm{risk})<\tau_c$，标记类别不确定并触发回退。$\tau_c$在验证集上按类别错误率与回退比例选择。最终评估前锁定配置，分布外及多标签请求单独分析。

\subsection{面向不同风险等级的编辑模块动态调度}

调度设计见表\ref{tab:ch3-dispatch}。低置信度高风险提示可触发预设安全响应；该回退属于工程机制，应分别报告回退前算法指标与包含回退的系统指标，防止将规则收益计入编辑贡献。

\begin{table}[htbp]
 \centering
 \caption{风险分级调度设计（阈值及效果待验证）}
 \label{tab:ch3-dispatch}
 \begin{tabular}{>{\raggedright\arraybackslash}p{0.16\textwidth}>{\raggedright\arraybackslash}p{0.25\textwidth}>{\raggedright\arraybackslash}p{0.43\textwidth}}
 \toprule
 输入状态 & 模块调用 & 目标行为及需检查问题\\ \midrule
 低风险 & 旁路，$a=0$ & 正常回答；检查风险漏检\\
 中风险 & Top-1，$a=\beta$ & 安全说明；检查部分调用是否仍有危险输出\\
 高风险 & Top-1，$a=1$ & 阻止危险指导；检查安全替代信息质量\\
 高风险且类别不确定 & 标记并进入安全回退 & 单独记录回退率与误拒率\\
 \bottomrule
 \end{tabular}
\end{table}

日志记录模型及模块版本、风险分数、类别分布、门控等级和回退标记。请求结束后清除当前路由状态；并发时采用请求局部状态或串行队列，防止模块跨请求误用。

\subsection{类别解耦与动态路由联合优化}

令$\Phi=\{A_c,B_c\}_{c=1}^{C}$，$\psi$为分类头参数，设计目标写为
\begin{equation}
 \begin{aligned}
 \mathcal J_3(\Phi,\psi)={}&\mathcal L_{\mathrm{edit}}(\Phi;\psi)
 +\lambda_{\mathrm{loc}}\mathcal R_{\mathrm{KL}}^{(3)}\\
 &+\lambda_{\mathrm{dec}}\mathcal R_{\mathrm{dec}}
 +\lambda_{\mathrm{route}}\mathcal L_{\mathrm{route}}(\psi).
 \end{aligned}
 \label{eq:ch3-joint}
\end{equation}
其中$\mathcal L_{\mathrm{edit}}$为风险输入经模块调用后的安全响应损失。$\mathcal R_{\mathrm{KL}}^{(3)}$在无害训练输入上比较单独激活各模块时的输出分布与基座输出分布，并对模块及样本取均值。采用全强度激活计算该项，使KL能够约束误激活带来的输出变化。

采用分阶段优化。先在训练集监督训练分类头，冻结分类头后在验证集确定门控阈值。逐类训练模块时使用标注类别，分别按标注等级以$\beta$和$1$缩放中、高风险样本的模块增量。最后固定路由配置，使用训练数据的预测门控少量校准模块。只使用预测类别与标注类别一致且实际激活模块的样本更新安全响应损失，路由错误另行统计。式\eqref{eq:ch3-joint}给出各阶段的优化项，梯度在硬Top-1选择处停止。主要测试使用预测路由，以标注类别替换预测类别的结果用于诊断类别识别误差。

\begin{enumerate}
 \item \textbf{输入：}原始模型、类别及等级训练集、无害训练集与验证集。
 \item \textbf{位置确定：}依据训练样本差异确定并固定编辑层及投影，记录汇总规则。
 \item \textbf{路由训练：}缓存冻结提示表示，训练风险和类别头，校准阈值。
 \item \textbf{类别编辑：}逐类训练低秩参数，使用非目标功能约束与无害保持项，冻结基座及其他模块。
\item \textbf{调用校准：}固定路由配置，在训练集上按预测路由校准模块，依据验证集选择停止点并冻结参数。
 \item \textbf{输出：}模块库、分类头、阈值及版本信息，在独立测试集评估。
\end{enumerate}

\section{实验结果与分析}

本节为执行与填写框架。所有数值、结果图和效果判断\textbf{待完成}，不将方法目标写为实测结论。

\subsection{实验设置}

与第二章共享主模型、基准实现版本、无害效用协议及生成配置；优先采用Llama-2-7B-Chat，最终以资源和兼容性核定。主实验选择四类风险，按原始问题和攻击模板去重划分；另建风险相邻正常请求集。记录三级标注规则、样本量、评估器及人工复核比例。

设置原始模型、DINM顺序编辑、共享低秩编辑、类别模块固定强度调用及本文完整方法作为主要对照。DINM逐类编辑用于干扰诊断，顺序编辑用于多类别主对比。共享低秩编辑采用同一位置和数据，使用相同风险门控，所有已激活请求调用同一个模块。在代码与模型兼容且资源允许时，增加采用同一评估协议的ToxEdit对照。

固定强度与分级强度使用同一组已训练模块权重、同一分类头与阈值，各请求的预测类别、低风险旁路及回退规则均保持一致。两组仅改变预测为中风险且实际激活模块的请求的强度，固定强度组取$1$，分级强度组取$\beta$。评估算法本身时，两组均关闭安全回退。评估系统时，两组使用相同回退规则，并分别报告两种口径的结果。

安全评估参考HarmBench对统一评估协议的设计，并参考SORRY-Bench按主题组织风险请求、对照人工标签检查自动判定的做法\cite{mazeika2024harmbench,xie2025sorrybench}。自动评估分别记录输入风险、输出危害和是否拒答，参照WildGuard对三类判定任务的划分\cite{han2024wildguard}。主实验的数据范围和评估器版本待实现后记录。

记录各组样本访问量、参数量、编辑时长及推理开销。多模块与共享LoRA总参数难以完全匹配时，补充等参数预算或等训练预算说明。主配置拟采用三个随机种子，实际数量据实记录。实验命令、配置和逐样本记录\textbf{待补充}。

\subsection{多类别编辑效果对比}

按类别报告ASR、SRR及宏平均，区分已见与未见改写，同时报告效用和ORR。表\ref{tab:ch3-main-results}仅为模板，不填写估计值或他人论文数值。
\begin{table}[htbp]
 \centering
 \caption{多类别安全编辑主结果模板（实验待完成）}
 \label{tab:ch3-main-results}
 \begin{tabular}{p{0.33\textwidth}cccc}
 \toprule
 方法 & 宏ASR$\downarrow$ & 宏SRR$\uparrow$ & 效用$\uparrow$ & ORR$\downarrow$\\ \midrule
 原始模型 & 待测 & 待测 & 待测 & 待测\\
 DINM & 待测 & 待测 & 待测 & 待测\\
 共享低秩编辑 & 待测 & 待测 & 待测 & 待测\\
 类别模块（固定强度） & 待测 & 待测 & 待测 & 待测\\
 本文完整方法 & 待测 & 待测 & 待测 & 待测\\
 \bottomrule
 \end{tabular}
\end{table}

分析按类别说明样本量、失败类型及安全与效用的变化。分类结果用于界定方法适用的风险类别。

\subsection{类别隔离与干扰分析}

依据式\eqref{eq:ch3-interference}绘制DINM、无解耦约束模块和带解耦约束模块的类别干扰矩阵。各方法使用相同的类别测试集。两组类别模块均固定以单位强度调用，关闭预测路由、旁路和回退。在目标类SRR相近时，比较非对角项的安全退化及无害请求的ORR。若目标类SRR差异较大，则同时报告目标类与非目标类变化，分析干扰减弱是否伴随目标编辑效果下降。表\ref{tab:ch3-interference-results}给出单个方法的填写格式，数值使用SRR百分点差。

另在预测路由下按同一类别顺序加入编辑，记录每次新增编辑后已有类别的SRR和无害请求的ORR。该结果用于观察实际调用下的变化，与固定调用的干扰矩阵分开解释。

\begin{table}[htbp]
 \centering
 \caption{逐类编辑的类别干扰矩阵（待测）}
 \label{tab:ch3-interference-results}
 \begin{tabular}{lcccc}
 \toprule
 编辑类别 & 测试类别1 & 测试类别2 & 测试类别3 & 测试类别4\\ \midrule
 类别1 & 待测 & 待测 & 待测 & 待测\\
 类别2 & 待测 & 待测 & 待测 & 待测\\
 类别3 & 待测 & 待测 & 待测 & 待测\\
 类别4 & 待测 & 待测 & 待测 & 待测\\
 \bottomrule
 \end{tabular}
\end{table}

\subsection{风险识别与分级性能分析}

报告风险召回率、精确率、类别宏F1、三级混淆矩阵和高风险漏检率。诊断时保持风险强度与回退规则固定，仅将预测类别替换为标注类别，记录响应指标变化。真实类别诊断与部署结果分别记录。标注一致性、校准过程和误判样例\textbf{待完成}。

\subsection{动态路由性能分析}

比较共享模块、固定强度Top-1及三级强度Top-1，记录旁路率、调用分布、回退率、路由额外延迟、首token延迟与吞吐量。比较响应效果时，固定强度与分级强度两组复用同一份预测路由输出。测量部署开销时，计入完整生成路径的路由计算时间。各组保持硬件、生成设置、随机种子、批大小和生成长度一致。Top-2若实施则单独列示。效果与开销\textbf{待测}。

\subsection{安全效用平衡与过度拒答分析}

在一般无害请求和风险相邻正常请求上分别报告ORR、有用性和任务效用。按附录B的判定规则，在各标注等级内计算处置符合率，即符合该等级要求的输出数占该等级样本数的比例。空响应和生成失败计为处置不符合。分级强度的主要对比同时报告标注中风险请求的处置符合率、风险相邻正常请求的ORR，以及标注高风险集合上的绝对ASR。该集合保留被误判为中、低风险的请求，按式\eqref{eq:asr}计算ASR，并另报无效响应比例。

在验证阶段预先确定可接受的高风险ASR差异范围，并在测试前冻结配置。若两组的高风险ASR落在预设差异范围内，比较ORR与中风险处置符合率。若超出范围，则如实报告安全性与误拒之间的权衡。风险相邻正常请求另按是否被判为中风险且激活模块分组，并报告激活子集的样本数和ORR。该子集直接受到调用强度变化的影响，低风险旁路在两组中一致。表\ref{tab:ch3-graded-results}记录关闭回退时的成对比较，启用回退后的系统结果另报。激活子集为空时，其ORR记为不适用。各等级得分及脱敏案例\textbf{待完成}。

\begin{table}[htbp]
 \centering
 \caption{固定与分级强度的处置效果（关闭回退，待测）}
 \label{tab:ch3-graded-results}
 \setlength{\tabcolsep}{4pt}
 \begin{tabular}{@{}lcccc@{}}
 \toprule
 调用方式 & \shortstack{标注高风险\\ASR} & \shortstack{标注中风险\\处置符合率} &
 \shortstack{风险相邻正常请求\\ORR} & \shortstack{其中激活子集\\ORR}\\ \midrule
 固定强度 & 待测 & 待测 & 待测 & 待测\\
 分级强度 & 待测 & 待测 & 待测 & 待测\\
 \bottomrule
 \end{tabular}
\end{table}

\subsection{消融与敏感性分析}

主消融分别移除非目标功能约束（$\lambda_{\mathrm{dec}}=0$）、将已激活的预测中风险请求强度设为$1$、以共享模块代替类别拆分。各组沿用相同的编辑位置、数据和评估协议，涉及重新训练的变体控制训练预算。固定强度组复用完整方法的模块及路由配置。敏感性实验采用单因素比较，重点考察秩$k$、解耦权重$\lambda_{\mathrm{dec}}$及风险阈值。消融表与曲线\textbf{待完成}。

\section{本章小结}

本章用非目标功能约束控制各类别模块的作用范围，并以预测类别选择模块、以预测等级确定调用强度。固定调用的干扰矩阵用于检验模块作用范围。端到端实验在相近高风险防护水平下比较固定和分级强度的中风险处置与相邻无害请求误拒情况。当前完成方法与实验框架，数值结果和效果结论待完成。第三章在原始基座上独立验证，与第二章模型的组合需另作兼容性评估。
